	\section{Planteamiento del Problema}
	Este capítulo presenta el análisis de las principales problemáticas identificadas en los procesos
	actuales de atención y gestión de quejas dentro de la DDP. Se describen las causas que originan
	dichas limitaciones operativas y tecnológicas, así como la propuesta de solución planteada para
	optimizar la administración de expedientes, mejorar la trazabilidad de la información y fortalecer
	la eficiencia institucional mediante una plataforma digital integral.
	
	\subsection{Problema general}
	La DDP enfrenta diversos retos operativos derivados de la gestión manual y parcialmente digitalizada
	de las quejas que presenta la comunidad politécnica. Hoy en día, buena parte del registro, la
	validación, el seguimiento y la resolución de expedientes descansa en mecanismos administrativos
	fragmentados, lo que se traduce en retrasos en la atención, dificultades de seguimiento y una
	capacidad limitada de supervisión institucional \cite{ipn_manual_procedimientos_ddp}.
	
	A ello se suma que la ausencia de una plataforma centralizada entorpece la comunicación entre las
	áreas que intervienen en el procedimiento y provoca inconsistencias en el flujo operativo, duplicidad
	de información y complicaciones para dar seguimiento oportuno a cada expediente. Esta situación
	repercute de forma directa en la eficiencia institucional, en la experiencia del usuario y en la
	capacidad de respuesta de la DDP ante situaciones que exigen atención inmediata.
	
	La administración de expedientes y evidencias digitales presenta además limitaciones importantes de
	organización, consulta y resguardo seguro de la información, lo cual eleva el riesgo de pérdida
	documental y de errores administrativos, y dificulta garantizar la integridad de los datos durante
	todo el ciclo de vida de la queja.
	
	Por último, la determinación preliminar de competencia y procedencia depende en gran medida de la
	revisión manual de las narrativas que redactan los usuarios, lo que incrementa la carga operativa del
	personal y puede generar inconsistencias en la clasificación inicial de los casos
	\cite{ipn_acuerdo622}. En consecuencia, la DDP encuentra dificultades para reducir sus tiempos de
	respuesta, abatir el rezago administrativo y sostener un flujo de atención eficiente y transparente.
	
	\subsection{Problemas específicos}
	Del análisis del flujo operativo institucional y de los procesos actuales de atención de quejas se
	desprenden los siguientes problemas específicos:
	
	\begin{itemize}
		\item \textbf{Ausencia de una plataforma centralizada:} no existe un punto único para la gestión
		integral de expedientes, evidencias y seguimiento de casos.
	
		\item \textbf{Dependencia de procesos manuales:} el registro, la validación y la canalización de
		las quejas se realizan sin apoyo de herramientas automatizadas.
	
		\item \textbf{Determinación de competencia lenta e inconsistente:} resulta difícil establecer con
		rapidez y bajo un criterio uniforme la competencia institucional de cada caso.
	
		\item \textbf{Trazabilidad escasa:} el ciclo de vida de las quejas deja poco rastro documental y el
		estado actual de los expedientes tiene visibilidad limitada.
	
		\item \textbf{Retrasos administrativos:} la revisión manual de documentos y evidencias digitales
		alarga los tiempos de atención.
	
		\item \textbf{Comunicación deficiente entre áreas:} las distintas áreas que participan en el
		proceso de atención carecen de canales integrados de coordinación.
	
		\item \textbf{Riesgos en el manejo de información sensible:} las evidencias digitales se
		administran sin mecanismos robustos de control de acceso y auditoría.
	
		\item \textbf{Falta de validación automatizada:} no existen mecanismos que detecten información
		incompleta, inconsistente o capturada de forma incorrecta por los usuarios.
	
		\item \textbf{Historial poco estructurado:} resulta complicado mantener el registro ordenado de
		actuaciones, acuerdos, notificaciones y resoluciones de cada expediente.
	
		\item \textbf{Limitaciones para la explotación de datos:} la infraestructura actual no permite
		generar estadísticas, indicadores ni reportes institucionales.
	\end{itemize}
	
	\subsection{Causas}
	Las problemáticas identificadas se originan en factores tecnológicos, operativos y organizacionales
	que inciden de manera directa en el desempeño de los procesos internos de la DDP.
	
	En primer lugar, una parte importante del procedimiento administrativo recae en actividades manuales
	del personal institucional: validación de información, clasificación de quejas, revisión documental
	y seguimiento de expedientes. Esa dependencia eleva la probabilidad de errores humanos,
	inconsistencias operativas y retrasos en la atención.
	
	En segundo lugar, la inexistencia de una arquitectura tecnológica centralizada fragmenta la
	información entre las distintas áreas y dificulta integrar los procesos institucionales de forma
	eficiente. Como resultado, los datos de una misma queja pueden encontrarse dispersos, duplicados o
	desactualizados, lo que afecta la continuidad operativa y la trazabilidad del expediente.
	
	Un tercer factor es la ausencia de mecanismos automatizados para analizar y clasificar las narrativas
	que presentan los usuarios. La determinación preliminar de competencia descansa hoy en la
	interpretación humana de los hechos descritos en cada queja, situación que incrementa la carga
	operativa y abre espacio a diferencias de criterio entre operadores.
	
	A lo anterior se añade la falta de herramientas especializadas para el control de accesos, la
	administración de evidencias digitales y el monitoreo de actividades, lo que limita la capacidad
	institucional de garantizar seguridad, privacidad e integridad documental.
	
	Por último, el crecimiento del volumen de información y la complejidad de los procesos
	administrativos evidencian las limitaciones de la infraestructura tecnológica tradicional, que
	dificulta la escalabilidad, la mantenibilidad y la modernización de los servicios digitales
	institucionales.
	
	\subsection{Alcances}
	El proyecto delimita su alcance a la Defensoría de los Derechos Politécnicos del Instituto
	Politécnico Nacional y abarca únicamente los procesos internos de gestión de quejas de dicha
	instancia. La plataforma será accesible desde navegador web en equipos de escritorio y portátiles;
	no se contempla el desarrollo de una aplicación móvil nativa.
	
	El sistema cubre el ciclo operativo completo de una queja, desde su recepción hasta su resolución, e
	incluye módulos diferenciados para cada rol institucional: Quejoso, Recepcionista, Área de Primer
	Contacto, Defensor, Subdefensor y Titular. Se integra además un módulo de clasificación automática
	asistida, basado en Procesamiento de Lenguaje Natural, cuya función es sugerir la competencia
	preliminar de cada queja \cite{reimers2019sbert}. En la etapa que documenta este trabajo dicho
	módulo opera sobre un corpus sintético construido por el equipo de desarrollo, por lo que sus
	resultados tienen carácter preliminar y quedan sujetos a la recalibración señalada en la
	Sección~\ref{sec:limitaciones}. El quejoso podrá consultar el estado de su caso en cualquier
	momento mediante un folio digital único, generado de forma automática al momento del registro.
	
	Quedan excluidos de forma explícita la integración con sistemas externos al IPN, la gestión de
	procesos distintos al flujo de quejas de la DDP y la emisión automatizada de resoluciones o
	dictámenes legales.
	
	\subsection{Limitaciones}
	\label{sec:limitaciones}
	La ejecución del proyecto está condicionada por los siguientes factores:
	
	\begin{itemize}
		\item \textbf{Tiempo:} al tratarse de un proyecto académico con entregas parciales, no todos los
		módulos institucionales se encuentran operativos en esta etapa. Los componentes desarrollados hasta
		el momento son el microservicio de gestión de quejas y el módulo clasificador de competencia; los
		restantes quedan para una entrega posterior.
	
		\item \textbf{Datos:} el entrenamiento y la validación del clasificador se realizaron con un
		corpus sintético de quejas, construido por el equipo de desarrollo. A la fecha de elaboración de
		este documento, la DDP aún no ha entregado los casos históricos reales y anonimizados
		solicitados, cuya gestión continúa en trámite. Dado que la calidad y el volumen de los datos de
		entrenamiento inciden directamente en la precisión del modelo, las métricas reportadas en este
		trabajo deben considerarse preliminares y quedan sujetas a recalibración una vez que se disponga
		de los expedientes reales.
	
		\item \textbf{Adopción institucional:} la efectividad del sistema depende de la disposición del
		personal de la DDP para incorporar la plataforma a sus procesos cotidianos y validar las
		clasificaciones que sugiere el modelo de inteligencia artificial.
	
		\item \textbf{Restricción ética y normativa:} el sistema no sustituye en ningún caso el criterio
		jurídico humano. Las resoluciones finales son responsabilidad exclusiva del personal de la
		Defensoría, conforme a las reglas de negocio establecidas y a la normativa aplicable en materia de
		protección de datos personales.
	\end{itemize}
